\documentclass[11pt,a4paper]{amsart}

\usepackage[T1]{fontenc}
\usepackage{lmodern}
\usepackage{amsmath,amssymb,mathtools}
\usepackage{microtype}
\usepackage{booktabs}
\usepackage{enumitem}
\usepackage[a4paper,margin=30mm]{geometry}
\usepackage[dvipsnames]{xcolor}
\usepackage[colorlinks=true,linkcolor=MidnightBlue,citecolor=MidnightBlue,urlcolor=MidnightBlue]{hyperref}
\hypersetup{
  pdftitle={Selberg Feedback and Approximate Divisor Data Do Not Stably Determine Square-Root Prime Error},
  pdfauthor={Mayk Loide Baccaro}
}

\newtheorem{theorem}{Theorem}[section]
\newtheorem{corollary}[theorem]{Corollary}
\newtheorem{proposition}[theorem]{Proposition}
\newtheorem{lemma}[theorem]{Lemma}
\theoremstyle{definition}
\newtheorem{definition}[theorem]{Definition}
\newtheorem{remark}[theorem]{Remark}

\DeclareMathOperator{\suppop}{supp}
\newcommand{\N}{\mathbb N}
\newcommand{\1}{\mathbf 1}
\newcommand{\Lambdaone}{\Lambda}

\title[Selberg feedback and prime error]{Selberg Feedback and Approximate Divisor Data\\Do Not Stably Determine Square-Root Prime Error}
\author[Mayk Loide Baccaro]{Mayk Loide Baccaro\\
Independent researcher\\
\href{mailto:m.baccaro7663@student.nu.edu}{\texttt{m.baccaro7663@student.nu.edu}}}
\date{18 August 2026}
\subjclass[2020]{Primary 11N05; Secondary 11M26, 11N37, 11A25}
\keywords{Selberg symmetry formula, von Mangoldt function, Dirichlet convolution, stability, inverse problem, Riemann hypothesis}
\begin{document}

\begin{abstract}
Selberg's symmetry formula is strong enough to prove the prime number theorem.
Approximate versions of its feedback can lose square-root-scale information
about prime error.  We construct positive weights on the actual
prime powers that are constant along every prime-power chain and satisfy the
Selberg symmetry formula with its classical $O(x)$ remainder, yet whose
summatory functions separate by $x^\rho$ for any prescribed
$\rho\in(1/2,1)$.  After scale normalization, each complete divisor-residual
vector through $2X$ has norm at most a fixed tolerance $\tau$ (and the two
vectors are therefore at distance at most $2\tau$), while the summatory errors
at $X$ differ by $\gg_\rho X^\rho/\log X$.  Thus every exact recovery map on
this class has condition number at least $\gg_\rho X^\rho/\log X$ in the
supremum norm.  The same construction has centered divisor residual, bounded
mean square, a compactly supported dyadic transform obstruction, and
controlled generalized von Mangoldt recursions at every order.  These results
isolate a precise information barrier for the stated approximate summaries.
Exact symbolic use of $\Lambda*\1=\log$ lies outside this approximation
problem, and the results leave the Riemann hypothesis unchanged.
\end{abstract}

\maketitle

\section{Introduction}

Write
\[
  \psi(x)=\sum_{n\le x}\Lambda(n),
  \qquad (f*g)(n)=\sum_{d\mid n}f(d)g(n/d),
\]
where $\Lambda(p^k)=\log p$ and $\Lambda(n)=0$ off the prime powers.  The identity
\begin{equation}\label{eq:exact-inversion}
  \Lambda*\1=\log
\end{equation}
determines $\Lambda$ exactly by M\"obius inversion.  Stable recovery from
approximate divisor data is a separate question.  This paper quantifies its
conditioning inside a narrow arithmetic class that preserves the most
conspicuous structure of $\Lambda$: nonnegativity, support on the actual prime
powers, and constancy on each chain $p,p^2,p^3,\ldots$.

The Riemann hypothesis is equivalent to a square-root-scale estimate for $\psi(x)-x$, with the usual logarithmic allowance; this formulation goes back to von Koch \cite{vonKoch1901}.  Selberg's elementary proof of the prime number theorem derives the symmetry formula
\begin{equation}\label{eq:selberg-classical}
 \sum_{n\le x}\Lambda(n)\log n
 +\sum_{mn\le x}\Lambda(m)\Lambda(n)
 =2x\log x+O(x)
\end{equation}
and closes a qualitative feedback argument \cite{Selberg1949}.  We test
whether positivity, the $O(x)$ Selberg remainder, and small error in
\eqref{eq:exact-inversion} force a much sharper summatory estimate.  The
constructed family gives a quantitative failure.

For a weight $a$ define
\begin{align}
 A_a(x)&=\sum_{n\le x}a(n),
 &E_a(x)&=A_a(x)-x,\label{eq:A-E}\\
 R_a(n)&=(a*\1)(n)-\log n,\label{eq:residual}\\
 S_a(x)&=\sum_{n\le x}a(n)\log n
       +\sum_{mn\le x}a(m)a(n),\label{eq:Selberg-stat}
\end{align}
where the second sum in \eqref{eq:Selberg-stat} is over ordered pairs.  Our main theorem constructs, at every large scale $X$, two admissible weights whose residual vectors $(R_a(n))_{n\le2X}$ are uniformly close but whose outputs $E_a(X)$ are far apart.

Section~\ref{sec:main} states the conditioning theorem and its square-root impossibility corollary.  Sections~\ref{sec:construction}--\ref{sec:selberg} build the fixed two-sign family and prove its arithmetic and Selberg properties.  Section~\ref{sec:compact} gives a compact dyadic transform consequence, and Section~\ref{sec:hierarchy} treats the generalized von Mangoldt recursion.  Appendix~\ref{app:dense} records a simpler dense model that exposes the underlying cancellation without the prime-power constraints.

Related results fall into three groups.  Beurling's generalized primes show how prime-number theorems respond to changes in the underlying prime system \cite{Beurling1937}; here the ordinary primes remain fixed and only the weight changes.  Quantitative remainder theorems for generalized primes require explicit analytic hypotheses \cite{DebruyneVindas2017}.  Ivi\'c introduced functions generalizing von Mangoldt's function, and Calder\'on Garc\'ia--Z\'arate developed a Selberg-type asymptotic variant for such functions \cite{Ivic1975,CalderonGarciaZarate1990}.  Generalized von Mangoldt sums also encode criteria equivalent to the Riemann hypothesis \cite{Balazard2008,BanksSinha2025}, while prime-sensitive Lindel\"of formulations give another exact equivalence \cite{GonekGrahamLee2020}.  None of these results supplies the uniform inverse estimate tested here.  Our theorem addresses conditioning of recovery from approximate arithmetic summaries on a fixed positive prime-power cone.

\section{The conditioning obstruction}\label{sec:main}

\begin{definition}\label{def:cone}
Let $\mathcal C$ be the set of weights $a:\N\to[0,\infty)$ such that
\[
  \suppop(a)\subseteq\{p^k:p\text{ prime},\ k\ge1\}
  \quad\text{and}\quad
  a(p^k)=a(p)\quad(k\ge1).
\]
For $X\ge2$, write
\[
  \mathcal R_X(a)=(R_a(n))_{1\le n\le\lfloor2X\rfloor}
\]
and equip residual vectors with the supremum norm.
\end{definition}

\begin{theorem}[Near-linear conditioning lower bound]\label{thm:conditioning}
Fix $\rho\in(1/2,1)$.  There are a function $\delta_\rho:\N\to\mathbb R$, a constant $C_\rho>0$, and $X_\rho\ge2$ with the following property.  For every $X\ge X_\rho$ and $0<\tau\le1$, the weights
\begin{equation}\label{eq:scaled-pair}
  a_{X,\pm}=\Lambda\pm\eta_X\delta_\rho,
  \qquad
  \eta_X=\frac{\tau}{C_\rho\log(2X)},
\end{equation}
belong to $\mathcal C$ and satisfy
\begin{align}
 \|\mathcal R_X(a_{X,\pm})\|_\infty&\le\tau,\label{eq:input-bound}\\
 |E_{a_{X,+}}(X)-E_{a_{X,-}}(X)|
 &\ge \frac{\tau}{\rho C_\rho}\frac{X^\rho}{\log(2X)}.\label{eq:output-separation}
\end{align}
Each $a_{X,\pm}$ also obeys
\begin{equation}\label{eq:scaled-selberg}
  S_{a_{X,\pm}}(x)=2x\log x+O_\rho(x)
  \qquad(x\ge2),
\end{equation}
with a constant independent of $X$ and $\tau$.

Consequently, if a map
\[
 \Phi_X:\mathcal R_X(\mathcal C)\longrightarrow\mathbb R
\]
recovers $E_a(X)$ exactly for every $a\in\mathcal C$, then its Lipschitz constant in the supremum norm satisfies
\begin{equation}\label{eq:lipschitz-lower}
  \operatorname{Lip}(\Phi_X)
  \ge \frac{1}{2\rho C_\rho}\frac{X^\rho}{\log(2X)}
  \qquad(X\ge X_\rho).
\end{equation}
In particular, for every $\delta\in(0,1/2)$ one can take $\rho=1-\delta$ and obtain
\[
  \operatorname{Lip}(\Phi_X)\gg_\delta
  \frac{X^{1-\delta}}{\log X}.
\]
\end{theorem}

The scale dependence in \eqref{eq:scaled-pair} is essential and explicit: the theorem measures the condition number of a family of finite-scale inverse problems.  The fixed perturbations from which the family is obtained appear in Theorem~\ref{thm:fixed-pair}.

\begin{corollary}[No uniform square-root correction]\label{cor:sqrt}
Fix $\varepsilon\in[0,1/2)$.  There is no constant $C_\varepsilon$ such that every $a\in\mathcal C$ and every $X\ge2$ satisfy
\begin{equation}\label{eq:forbidden-bound}
 |A_a(X)-\psi(X)|
 \le C_\varepsilon X^{1/2+\varepsilon}
 \max_{n\le2X}|R_a(n)|.
\end{equation}
The conclusion remains true after restricting to weights that satisfy \eqref{eq:scaled-selberg} with a uniform $O_\rho(x)$ constant.
\end{corollary}

\begin{proof}
Choose $\rho\in(1/2+\varepsilon,1)$ and use one of the weights in \eqref{eq:scaled-pair}.  Since
\[
 A_{a_{X,\pm}}(X)-\psi(X)
 =\pm\eta_X\sum_{n\le X}\delta_\rho(n),
\]
the asymptotic established in Proposition~\ref{prop:fixed-sums} gives, for both signs and all sufficiently large $X$,
\[
 |A_{a_{X,\pm}}(X)-\psi(X)|
 \ge \frac{\tau}{2\rho C_\rho}
       \frac{X^\rho}{\log(2X)}.
\]
The right side of \eqref{eq:forbidden-bound} is at most
$C_\varepsilon\tau X^{1/2+\varepsilon}$.  Their ratio tends to infinity.
\end{proof}

\begin{table}[t]
\centering
\caption{Properties of the fixed and scale-normalized pairs.  All constants may depend on $\rho$.}
\label{tab:comparison}
\small
\begin{tabular}{@{}p{0.39\textwidth}cc@{}}
\toprule
Property & Fixed pair & Normalized pair at $X$\\
\midrule
Actual prime-power support & yes & yes\\
Strictly positive at every prime power & yes & yes\\
Constant on every $p^k$ chain & yes & yes\\
Centered divisor mean & $O(x^\rho)$ & $O(x^\rho/\log X)$\\
Divisor residual mean square & $O(x)$ & $O(x/\log^2 X)$\\
Complete residual through $2X$ & $O(\log X)$ & at most $\tau$\\
Selberg statistic & $2x\log x+O(x)$ & uniform $2x\log x+O(x)$\\
Two-sign output separation at $X$ & $\asymp X^\rho$ & $\gg \tau X^\rho/\log X$\\
\bottomrule
\end{tabular}
\end{table}

\section{A centered perturbation on the actual prime powers}\label{sec:construction}

Fix $\rho\in(1/2,1)$ for the remainder of Sections~\ref{sec:construction}--\ref{sec:compact}, and define
\begin{equation}\label{eq:c-p}
  c_2=-\sum_{\substack{p\text{ prime}\\p\text{ odd}}}
       \frac{p^{\rho-1}\log p}{p-1},
  \qquad
  c_p=p^{\rho-1}\log p\quad(p\text{ odd}).
\end{equation}
The series defining $c_2$ converges because its summand has order
$\log p/p^{2-\rho}$ and $2-\rho>1$.  Extend the coefficients to an arithmetic function by
\begin{equation}\label{eq:delta}
  \delta_\rho(p^k)=c_p\quad(k\ge1),
  \qquad
  \delta_\rho(n)=0\quad(n\text{ not a prime power}).
\end{equation}
The correction on the $2$-power chain enforces the exact centering identity
\begin{equation}\label{eq:centering}
  \sum_p\frac{c_p}{p-1}=0.
\end{equation}

Set
\begin{equation}\label{eq:constants}
 C_\rho=\max\left\{3^{\rho-1},\frac{|c_2|}{\log2}\right\},
 \qquad
 D_\rho=\sum_{n\ge1}\frac{|\delta_\rho(n)|}{n}
 =\sum_p\frac{|c_p|}{p-1}.
\end{equation}
For odd primes, $|c_p|/\log p=p^{\rho-1}\le3^{\rho-1}$, so
$|c_p|\le C_\rho\log p$ for every prime $p$.  Both constants are finite,
and $C_\rho>0$.

\begin{proposition}[Size of the perturbation]\label{prop:fixed-sums}
As $x\to\infty$,
\begin{align}
 \sum_{n\le x}|\delta_\rho(n)|&=O_\rho(x^\rho),\label{eq:U}\\
 H_\rho(x):=\sum_{n\le x}\delta_\rho(n)
 &=\frac{x^\rho}{\rho}+o_\rho(x^\rho).\label{eq:H}
\end{align}
\end{proposition}

\begin{proof}
Chebyshev's estimate $\vartheta(y)=\sum_{p\le y}\log p=O(y)$ and partial summation give
\[
 \sum_{p\le y}p^{\rho-1}\log p=O_\rho(y^\rho).
\]
The levels $p^k$ with $k\ge2$ contribute $O_\rho(x^{\rho/2}\log x)$, and the correction along the powers of $2$ contributes $O_\rho(\log x)$.  This proves \eqref{eq:U}.

For the signed sum, the prime level contributes
\[
 \sum_{p\le x}p^{\rho-1}\log p
 =\frac{x^\rho}{\rho}+o_\rho(x^\rho)
\]
by the prime number theorem and partial summation.  The higher prime-power levels and the $2$-chain correction are $o_\rho(x^\rho)$ by the bounds just given.  Hence \eqref{eq:H} follows.
\end{proof}

\begin{theorem}[Fixed two-sign family]\label{thm:fixed-pair}
For every
\begin{equation}\label{eq:eta-small}
 0<\eta\le \min\{1,(2C_\rho)^{-1}\},
\end{equation}
the weights
\[
 a_\pm=\Lambda\pm\eta\delta_\rho
\]
belong to $\mathcal C$, are strictly positive on every prime power, and satisfy
\begin{equation}\label{eq:fixed-separation}
 A_{a_+}(x)-A_{a_-}(x)
 =\frac{2\eta}{\rho}x^\rho+o_\rho(x^\rho).
\end{equation}
\end{theorem}

\begin{proof}
Support and chain constancy follow from \eqref{eq:delta}.  At $p^k$,
\[
 a_\pm(p^k)=\log p\pm\eta c_p
 \ge(1-\eta C_\rho)\log p\ge\tfrac12\log p>0.
\]
Equation \eqref{eq:fixed-separation} is twice $\eta$ times \eqref{eq:H}.
\end{proof}

\section{The divisor residual}\label{sec:divisor}

Let
\begin{equation}\label{eq:r-def}
 r_\rho=\delta_\rho*\1.
\end{equation}
Unique factorization turns the convolution into the explicit additive function
\begin{equation}\label{eq:r-explicit}
 r_\rho(n)=\sum_p c_pv_p(n),
\end{equation}
where $v_p(n)$ is the exponent of $p$ in $n$.  Thus
\begin{equation}\label{eq:r-pointwise}
 |r_\rho(n)|
 \le C_\rho\sum_pv_p(n)\log p
 =C_\rho\log n.
\end{equation}

\begin{proposition}[Centered mean and bounded mean square]\label{prop:residual-moments}
For $x\ge2$,
\begin{align}
 \sum_{n\le x}r_\rho(n)&=O_\rho(x^\rho),\label{eq:r-mean}\\
 \sum_{n\le x}|r_\rho(n)|^2&=O_\rho(x).\label{eq:r-L2}
\end{align}
Consequently, the residuals of the fixed weights satisfy
\[
 R_{a_\pm}(n)=\pm\eta r_\rho(n)
\]
with centered mean, bounded mean square, and the pointwise estimate
$|R_{a_\pm}(n)|\le\eta C_\rho\log n$.
\end{proposition}

\begin{proof}
For each prime $p\le x$,
\begin{equation}\label{eq:vp-first}
 \sum_{n\le x}v_p(n)
 =\sum_{\substack{k\ge1\\p^k\le x}}\left\lfloor\frac{x}{p^k}\right\rfloor
 =\frac{x}{p-1}+O\!\left(1+\frac{\log x}{\log p}\right).
\end{equation}
Insert \eqref{eq:vp-first} into \eqref{eq:r-explicit}.  The complete main term vanishes by \eqref{eq:centering}.  Partial summation with $\vartheta(y)=O(y)$ bounds the omitted tail by
\[
 x\sum_{p>x}\frac{|c_p|}{p-1}=O_\rho(x^\rho),
\]
while
\[
 \sum_{p\le x}|c_p|\left(1+\frac{\log x}{\log p}\right)
 =O_\rho(x^\rho).
\]
This proves \eqref{eq:r-mean}.

For the second moment, use
\begin{align}
 \sum_{n\le x}v_p(n)^2
 &\le x\sum_{k,l\ge1}p^{-\max(k,l)}
 =x\frac{p+1}{(p-1)^2},\label{eq:vp-square}\\
 \sum_{n\le x}v_p(n)v_q(n)
 &\le x\sum_{k,l\ge1}p^{-k}q^{-l}
 =\frac{x}{(p-1)(q-1)}\quad(p\ne q).\label{eq:vp-cross}
\end{align}
Expanding the square in \eqref{eq:r-explicit}, taking absolute values on the cross terms, and applying \eqref{eq:vp-square}--\eqref{eq:vp-cross} gives
\[
 \sum_{n\le x}|r_\rho(n)|^2
 \le x\sum_p|c_p|^2\frac{p+1}{(p-1)^2}
 +2x\sum_{p<q}\frac{|c_pc_q|}{(p-1)(q-1)}.
\]
The first series converges because its general term is
$O_\rho((\log p)^2/p^{3-2\rho})$; the second is bounded by $xD_\rho^2$.
\end{proof}

\section{Selberg feedback survives the perturbation}\label{sec:selberg}

\begin{proposition}[Uniform Selberg remainder]\label{prop:selberg}
Let $a=\Lambda+e$, where $e=\pm\eta\delta_\rho$ and $|\eta|\le\eta_0$ for a fixed $\eta_0$.  Then
\begin{equation}\label{eq:selberg-perturbed}
 S_a(x)=2x\log x+O_{\rho,\eta_0}(x).
\end{equation}
\end{proposition}

\begin{proof}
Expanding the ordered convolution gives
\begin{equation}\label{eq:S-difference}
 S_a(x)-S_\Lambda(x)
 =\sum_{n\le x}e(n)\log n
 +2\sum_{mn\le x}e(m)\Lambda(n)
 +\sum_{mn\le x}e(m)e(n).
\end{equation}
The first term is $O_{\rho,\eta_0}(x^\rho\log x)=O_{\rho,\eta_0}(x)$ by \eqref{eq:U}.  Choose $C_\psi$ so that $\psi(y)\le C_\psi y$ for $y\ge1$.  Then
\[
 \sum_{mn\le x}|e(m)|\Lambda(n)
 \le C_\psi x\sum_{m\ge1}\frac{|e(m)|}{m}
 \le C_\psi\eta_0D_\rho x.
\]
Also
\[
 \sum_{mn\le x}|e(m)e(n)|
 \le x\left(\sum_{m\ge1}\frac{|e(m)|}{m}\right)^2
 \le\eta_0^2D_\rho^2x.
\]
Combine these estimates with Selberg's formula \eqref{eq:selberg-classical}.
\end{proof}

\begin{proof}[Proof of Theorem~\ref{thm:conditioning}]
Use the perturbation \eqref{eq:delta}.  Since $X\ge2$ and $\tau\le1$,
\[
 \eta_XC_\rho=\frac{\tau}{\log(2X)}\le\frac1{\log4}<1.
\]
The positivity argument from Theorem~\ref{thm:fixed-pair} therefore applies globally to $a_{X,\pm}$.  Equations \eqref{eq:residual}, \eqref{eq:exact-inversion}, and \eqref{eq:r-pointwise} yield, for $n\le2X$,
\[
 |R_{a_{X,\pm}}(n)|
 =\eta_X|r_\rho(n)|
 \le\eta_XC_\rho\log(2X)=\tau.
\]
By \eqref{eq:H}, enlarge $X_\rho$ until
$|H_\rho(X)|\ge X^\rho/(2\rho)$.  Then
\[
 |E_{a_{X,+}}(X)-E_{a_{X,-}}(X)|
 =2\eta_X|H_\rho(X)|
 \ge\frac{\tau}{\rho C_\rho}\frac{X^\rho}{\log(2X)}.
\]
Proposition~\ref{prop:selberg} applies uniformly because
$\eta_X\le(C_\rho\log4)^{-1}$.  Finally,
\[
 \|\mathcal R_X(a_{X,+})-\mathcal R_X(a_{X,-})\|_\infty
 \le2\tau.
\]
Taking the quotient of the output separation by this input distance proves
\eqref{eq:lipschitz-lower}.
\end{proof}

\section{A compact dyadic transform}\label{sec:compact}

The full residual vector in Theorem~\ref{thm:conditioning} is not the only summary that loses the $x^\rho$ direction.  A simple dyadic combination produces an exactly compactly supported kernel.

Define
\begin{equation}\label{eq:K}
 K(u)=
 \begin{cases}
 0,&u\le0,\\
 \tfrac12,&0<u\le1,\\
 \dfrac{1}{2u^2},&u>1,
 \end{cases}
\end{equation}
and set
\begin{equation}\label{eq:transform}
 T_a(X)=\sum_{n\ge1}a(n)K(n/X),
 \qquad
 M(X)=\int_1^\infty K(t/X)\,dt,
 \qquad
 \mathcal E_a(X)=M(X)-T_a(X).
\end{equation}
The series converges absolutely for the weights considered here.

\begin{proposition}[Compact dyadic separation]\label{prop:compact}
For the fixed pair $a_\pm=\Lambda\pm\eta\delta_\rho$, define
\[
 \mathcal D_a(X)=4\mathcal E_a(X)-\mathcal E_a(2X).
\]
Then
\begin{equation}\label{eq:compact-separation}
 \mathcal D_{a_+}(X)-\mathcal D_{a_-}(X)
 =-2\eta(4-2^\rho)I_\rho X^\rho+o_\rho(X^\rho),
 \qquad
 I_\rho=\frac1{\rho(2-\rho)}.
\end{equation}
Equivalently, the perturbation is tested against the compactly supported kernel
\begin{equation}\label{eq:W}
 W(u)=4K(u)-K(u/2)=
 \begin{cases}
 \tfrac32,&0<u\le1,\\
 \dfrac{2}{u^2}-\tfrac12,&1<u\le2,\\
 0,&u>2.
 \end{cases}
\end{equation}
If
\[
 \mathcal Q_a(Y)=\sum_{\substack{j\ge0\\2^j\le Y}}
 |\mathcal D_a(2^j)|^2,
\]
then for every $\varepsilon<2\rho-1$, at least one of $a_+$ and $a_-$ fails the estimate
$\mathcal Q_a(Y)=O(Y^{1+\varepsilon})$.
\end{proposition}

\begin{proof}
Stieltjes summation and \eqref{eq:H} give
\begin{equation}\label{eq:mellin-asymptotic}
 \sum_{n\ge1}\delta_\rho(n)K(n/X)
 =I_\rho X^\rho+o_\rho(X^\rho),
\end{equation}
where
\[
 I_\rho=\int_0^\infty u^{\rho-1}K(u)\,du
 =\frac1{2\rho}+\frac1{2(2-\rho)}
 =\frac1{\rho(2-\rho)}.
\]
Since the common main term $M$ cancels between signs,
\[
 \mathcal D_{a_+}(X)-\mathcal D_{a_-}(X)
 =-2\eta\sum_{n\ge1}\delta_\rho(n)W(n/X).
\]
Equations \eqref{eq:mellin-asymptotic} and \eqref{eq:W} yield
\eqref{eq:compact-separation}.

For all sufficiently large dyadic $X$,
\[
 |\mathcal D_{a_+}(X)|^2+|\mathcal D_{a_-}(X)|^2
 \ge\frac12|\mathcal D_{a_+}(X)-\mathcal D_{a_-}(X)|^2
 \gg_{\rho,\eta}X^{2\rho}.
\]
Summing over dyadic $X\le Y$ gives
$\mathcal Q_{a_+}(Y)+\mathcal Q_{a_-}(Y)\gg_{\rho,\eta}Y^{2\rho}$.
At least one fixed sign therefore violates every $O(Y^{1+\varepsilon})$ bound with
$1+\varepsilon<2\rho$.
\end{proof}

\section{Generalized von Mangoldt recursions}\label{sec:hierarchy}

Following the generalized von Mangoldt hierarchy studied by Ivi\'c and in
subsequent Selberg-type variants \cite{Ivic1975,CalderonGarciaZarate1990,BanksSinha2025}, define
the functions by
\begin{equation}\label{eq:Lambda-k}
 \Lambda_k=\mu*\log^k.
\end{equation}
They are nonnegative and satisfy the recursion
\begin{equation}\label{eq:Lambda-recursion}
 \Lambda_{k+1}=\Lambda_k\log+\Lambda_k*\Lambda
 \qquad(k\ge1).
\end{equation}
For a positive weight $a$, define its companion hierarchy by
\begin{equation}\label{eq:G-recursion}
 G_{a,1}=a,
 \qquad
 G_{a,k+1}=G_{a,k}\log+G_{a,k}*a.
\end{equation}
This recursion agrees with \eqref{eq:Lambda-recursion} at $a=\Lambda$ and preserves nonnegativity.

\begin{lemma}[Weighted prefix bound]\label{lem:prefix}
Let $f,g\ge0$ satisfy
\[
 \sum_{n\le y}f(n)\le Qy\log^m(2y),
 \qquad
 \sum_{n\le y}g(n)\le Cy
 \quad(y\ge1).
\]
With $q=1+1/\log2$,
\[
 \sum_{n\le x}(f*g)(n)
 \le qCQx\log^{m+1}(2x).
\]
\end{lemma}

\begin{proof}
Partial summation gives
\[
 \sum_{n\le x}\frac{g(n)}n\le C(1+\log x)
 \le qC\log(2x).
\]
Therefore
\[
 \sum_{uv\le x}f(u)g(v)
 \le Qx\log^m(2x)\sum_{v\le x}\frac{g(v)}v,
\]
which is the asserted estimate.
\end{proof}

\begin{theorem}[Order-uniform hierarchy control]\label{thm:hierarchy}
Let $a=a_+$ or $a_-$ be a fixed weight from Theorem~\ref{thm:fixed-pair}.  There exist constants $B,C\ge1$, depending on $\rho$ and $\eta$ but not on $k$ or $x$, such that
\begin{equation}\label{eq:hierarchy-bound}
 \sum_{n\le x}|G_{a,k}(n)-\Lambda_k(n)|
 \le BC^{k-2}x\log^{k-2}(2x)
 \qquad(k\ge2,\ x\ge1).
\end{equation}
\end{theorem}

\begin{proof}
Write $e=a-\Lambda$ and
$\Delta_k=G_{a,k}-\Lambda_k$.  Choose constants $C_U,C_\psi$ such that
\[
 \sum_{n\le x}|e(n)|\le C_Ux^\rho,
 \qquad
 \psi(x)\le C_\psi x
 \quad(x\ge1),
\]
and retain $D=\sum|e(n)|/n<\infty$.  The expansion at order two is
\[
 \Delta_2=e\log+e*\Lambda+\Lambda*e+e*e.
\]
Since
$x^{\rho-1}\log x\le[e(1-\rho)]^{-1}$ for $x\ge1$, the argument of
Proposition~\ref{prop:selberg} gives
\begin{equation}\label{eq:base-B0}
 \sum_{n\le x}|\Delta_2(n)|\le B_0x,
 \quad
 B_0=\frac{C_U}{e(1-\rho)}+2C_\psi D+D^2.
\end{equation}

Set
\[
 q=1+\frac1{\log2},
 \qquad C_\Lambda=1+qC_\psi,
 \qquad C_a=C_\psi+C_U.
\]
Induction from \eqref{eq:Lambda-recursion} and Lemma~\ref{lem:prefix} yields
\begin{equation}\label{eq:Lambda-prefix}
 \sum_{n\le x}\Lambda_k(n)
 \le C_\psi C_\Lambda^{k-1}x\log^{k-1}(2x).
\end{equation}
Subtracting \eqref{eq:Lambda-recursion} from \eqref{eq:G-recursion} gives
\begin{equation}\label{eq:Delta-recursion}
 \Delta_{k+1}=\Delta_k\log+\Delta_k*a+\Lambda_k*e.
\end{equation}

Take
\[
 B\ge\max\{1,B_0,C_\psi D\},
 \qquad
 C\ge2\max\{1+qC_a,C_\Lambda\}.
\]
Assume \eqref{eq:hierarchy-bound} at order $k$.  The first two terms in
\eqref{eq:Delta-recursion}, using Lemma~\ref{lem:prefix} for the convolution, contribute at most
\[
 B C^{k-2}(1+qC_a)x\log^{k-1}(2x)
 \le\tfrac12BC^{k-1}x\log^{k-1}(2x).
\]
Using \eqref{eq:Lambda-prefix} and $\sum|e(n)|/n=D$, the last term contributes at most
\[
 C_\psi D C_\Lambda^{k-1}x\log^{k-1}(2x)
 \le\tfrac12BC^{k-1}x\log^{k-1}(2x).
\]
This closes the induction.
\end{proof}

Theorem~\ref{thm:hierarchy} is deliberately order-uniform: it prints the full
$C^{k-2}$ dependence, with all growth in $k$ visible.  For this constructed
pair, every displayed companion recursion obeys the stated order-uniform
envelope while the original $x^\rho$ summatory separation persists.  These
bounds alone therefore do not exclude the pair.

\section{Concluding perspective}

The construction separates four statements that can look deceptively close.
\begin{enumerate}[label=\textup{(\roman*)},leftmargin=2.2em]
 \item The exact convolution equation $a*\1=\log$ identifies $a=\Lambda$.
 \item A small residual $a*\1-\log$ permits unstable recovery of $A_a(X)-x$.
 \item Positivity, actual prime-power support, chain constancy, centered first and second residual moments, and an $O(x)$ Selberg remainder still permit poor square-root conditioning.
 \item Higher generalized von Mangoldt recursions stay quantitatively controlled along the same unstable family.
\end{enumerate}
Exact algebraic uniqueness coexists with an inverse-problem condition number
that grows almost linearly with the observation scale.  Any successful
square-root recovery argument from approximate divisor or Selberg data must
introduce a norm, cancellation principle, spectral constraint, or analytic
input that detects the signed prime-level direction constructed here.

\appendix

\section{A dense warm-up}\label{app:dense}

The following model removes the prime-power restrictions and displays the same Selberg cancellation in one line of construction.

\begin{proposition}\label{prop:dense}
Fix $c>0$ and $\rho\in(1/2,1)$, and set
\[
 a(n)=1+c\rho\,n^{\rho-1}.
\]
Then
\begin{align}
 \sum_{n\le x}a(n)&=x+cx^\rho+O_{c,\rho}(1),\label{eq:dense-sum}\\
 S_a(x)&=2x\log x+O_{c,\rho}(x).
\end{align}
\end{proposition}

\begin{proof}
Euler summation gives
$\sum_{n\le x}n^{\rho-1}=x^\rho/\rho+O_\rho(1)$, proving
\eqref{eq:dense-sum}.  For the constant weight $1$,
\[
 \sum_{n\le x}\log n+\sum_{mn\le x}1
 =2x\log x+O(x).
\]
If $h(n)=c\rho n^{\rho-1}$, then
\[
 \sum_{n\le x}h(n)\log n=O(x^\rho\log x)=O(x),
\]
and
\[
 \sum_{mn\le x}h(m)
 \le x\sum_{m\ge1}\frac{h(m)}m=O(x).
\]
The same convergent weighted sum bounds the quadratic term
$\sum_{mn\le x}h(m)h(n)$ by $O(x)$.  Expanding $S_{1+h}$ proves the result.
\end{proof}

\enlargethispage{\baselineskip}
\section{Verification remarks}\label{app:repro}

The construction uses no numerical fitting.  The definition of $c_2$ makes
the centering identity \eqref{eq:centering} immediate once the defining series
and $D_\rho$ are known to converge.  In Proposition~\ref{prop:fixed-sums}, only the prime level supplies
the term $x^\rho/\rho$; the higher prime powers and the corrected $2$-chain
are lower order.  The residual calculation \eqref{eq:r-explicit} uses the
diagonal density $(p+1)/(p-1)^2$ and the off-diagonal density
$1/((p-1)(q-1))$, while \eqref{eq:S-difference} uses the factor $2$ from the
ordered cross term.

For the compact transform, the untruncated Mellin coefficient is
$\int_0^\infty u^{\rho-1}K(u)\,du=1/(\rho(2-\rho))$.  The dyadic combination
$4K(u)-K(u/2)$ has compact support; $K$ itself has an unbounded tail.  The constants
$B$ and $C^{k-2}$ in \eqref{eq:hierarchy-bound} are kept explicit because
separate bounds for fixed $k$ would not yield order-uniform control.

\section*{Data and code availability}

The construction and every asymptotic estimate used in the theorems are
contained in the manuscript; no external data or numerical experiment is
required.  An accompanying Lean~4 development checks the finite conditioning
algebra: the distinguished correction, positivity of the signed
perturbations, residual scaling, two-sign separation, and the resulting
Lipschitz lower bound.  The analytic construction, asymptotic estimates, and
order-uniform hierarchy remain manuscript-level proofs.

\section*{Assistance disclosure}

AI systems were used extensively in mathematical derivation, Lean proof
development, literature discovery, organization, and typesetting.  The author
remains responsible for every mathematical statement, proof, citation, and
submission decision.

\bibliographystyle{amsplain}
\bibliography{references}

\end{document}
